\documentclass[12pt, a4paper]{article}
\usepackage[spanish]{babel}
\usepackage{geometry}
\usepackage{fancyhdr}
\usepackage{xcolor}
\usepackage{graphicx}
\usepackage{hyperref}

\geometry{margin=1in}

\pagestyle{fancy}
\fancyhf{}
\rhead{Ingeniería de Software}
\lhead{Metodologías de Manejo de Proyectos}
\cfoot{\thepage}

\title{\textbf{Metodologías de Manejo de Proyectos de Software}}
\author{Segundo Parcial - Ingeniería de Software}
\date{\today}

\begin{document}

\maketitle

\tableofcontents
\newpage

\section{Introducción}

Las metodologías de manejo de proyectos de software son un conjunto de procedimientos y prácticas que permiten organizar, planificar y controlar el desarrollo de proyectos informáticos de manera eficiente. Estas metodologías proporcionan marcos de referencia que facilitan la comunicación entre equipos, establecen tiempos de entrega, aseguran la calidad del código y optimizan el uso de recursos.

\section{Investigación Preliminar}

\subsection{Preguntas Clave}

Para comprender adecuadamente las metodologías de desarrollo de software, es fundamental responder a las siguientes preguntas:

\begin{itemize}
    \item ¿Cuántas metodologías existen?
    \item ¿Cuáles son las principales?
    \item ¿Qué es cada una?
    \item ¿Para qué sirve?
    \item ¿Cómo se implementa?
\end{itemize}

\section{Metodologías Principales}

\subsection{Agile}

\subsubsection{¿Qué es?}
Agile es una metodología ágil que enfatiza la iteración rápida, la adaptabilidad y la colaboración continua con el cliente. Se basa en ciclos cortos de desarrollo (sprints) de 1 a 4 semanas.

\subsubsection{¿Para qué sirve?}
Agile es ideal para proyectos donde los requisitos pueden cambiar frecuentemente, permitiendo flexibilidad y respuesta rápida a nuevas necesidades del cliente.

\subsubsection{¿Cómo se implementa?}
\begin{itemize}
    \item Reuniones diarias (stand-ups)
    \item División del trabajo en historias de usuario
    \item Sprints iterativos
    \item Retroalimentación continua del cliente
    \item Entrega incremental de funcionalidades
\end{itemize}

\subsection{Waterfall (Cascada)}

\subsubsection{¿Qué es?}
Waterfall es una metodología clásica y lineal donde cada fase del proyecto debe completarse antes de pasar a la siguiente (requisitos, diseño, implementación, pruebas, despliegue).

\subsubsection{¿Para qué sirve?}
Es adecuada para proyectos con requisitos bien definidos y estables, donde es importante documentación completa y planificación anticipada.

\subsubsection{¿Cómo se implementa?}
\begin{itemize}
    \item Análisis exhaustivo de requisitos
    \item Diseño detallado antes de la codificación
    \item Fases secuenciales y bien definidas
    \item Documentación extensa
    \item Pruebas al final del desarrollo
\end{itemize}

\subsection{Lean}

\subsubsection{¿Qué es?}
Lean es una metodología que se enfoca en eliminar desperdicios y optimizar procesos, buscando entregar máximo valor con mínimos recursos.

\subsubsection{¿Para qué sirve?}
Lean es útil para mejorar la eficiencia de proyectos existentes, reducir tiempos de ciclo y aumentar la productividad del equipo.

\subsubsection{¿Cómo se implementa?}
\begin{itemize}
    \item Identificación y eliminación de actividades sin valor
    \item Automatización de procesos repetitivos
    \item Mejora continua (Kaizen)
    \item Enfoque en el flujo de trabajo
    \item Reducción de tiempos de espera
\end{itemize}

\subsection{Black Belt (Lean Six Sigma)}

\subsubsection{¿Qué es?}
Black Belt es un nivel de certificación dentro de la metodología Six Sigma, enfocada en reducir variabilidad y defectos en procesos.

\subsubsection{¿Para qué sirve?}
Black Belt se utiliza para proyectos de mejora de procesos complejos, garantizando calidad mediante análisis estadístico riguroso.

\subsubsection{¿Cómo se implementa?}
\begin{itemize}
    \item Análisis estadístico detallado (DMAIC)
    \item Medición de variabilidad
    \item Implementación de controles de calidad
    \item Documentación de mejoras
    \item Capacitación del equipo
\end{itemize}

\section{Otras Metodologías Relevantes}

\subsection{DevOps}

\subsubsection{¿Qué es?}
DevOps es una metodología que integra el desarrollo (Dev) y las operaciones (Ops), promoviendo la colaboración, automatización y comunicación continua entre equipos para acelerar el ciclo de vida del software.

\subsubsection{¿Para qué sirve?}
DevOps es útil para reducir el tiempo entre desarrollo e implementación, mejorar la confiabilidad del sistema, automatizar despliegues y crear una cultura de responsabilidad compartida.

\subsubsection{¿Cómo se implementa?}
\begin{itemize}
    \item Integración continua (CI)
    \item Despliegue continuo (CD)
    \item Automatización de infraestructura (Infrastructure as Code)
    \item Monitoreo y logging centralizado
    \item Colaboración entre desarrolladores y operarios
    \item Pipelines automatizados de prueba y despliegue
\end{itemize}

\subsection{Scrum}

\subsubsection{¿Qué es?}
Scrum es un framework ágil que proporciona un conjunto de roles, eventos y artefactos para gestionar proyectos complejos mediante iteraciones cortas llamadas sprints.

\subsubsection{¿Para qué sirve?}
Scrum es ideal para equipos que necesitan estructura adicional en su proceso ágil, con roles claramente definidos y ceremoniasque facilitan la organización y el seguimiento del progreso.

\subsubsection{¿Cómo se implementa?}
\begin{itemize}
    \item Product Owner: define y prioriza requisitos
    \item Scrum Master: facilita el proceso y elimina obstáculos
    \item Development Team: implementa el producto
    \item Product Backlog: lista de requisitos priorizados
    \item Sprint Backlog: tareas del sprint actual
    \item Ceremonias: planificación de sprint, stand-up diario, revisión y retrospectiva
\end{itemize}

\subsection{Kanban}

\subsubsection{¿Qué es?}
Kanban es un sistema visual de gestión de trabajo que utiliza tableros (físicos o digitales) para representar el flujo de tareas, limitando el trabajo en progreso para mejorar la eficiencia.

\subsubsection{¿Para qué sirve?}
Kanban es excelente para equipos que requieren flexibilidad, visualización del trabajo, mejora continua y limitación de trabajo simultáneo para evitar cuellos de botella.

\subsubsection{¿Cómo se implementa?}
\begin{itemize}
    \item Crear tablero visual con columnas (Por hacer, En progreso, Hecho)
    \item Establecer límites de trabajo en progreso (WIP limits)
    \item Usar tarjetas para representar tareas
    \item Medir y optimizar el flujo
    \item Reuniones de sincronización regulares
    \item Mejora continua basada en métricas
\end{itemize}

\subsection{Extreme Programming (XP)}

\subsubsection{¿Qué es?}
Extreme Programming es una metodología ágil que enfatiza prácticas técnicas rigurosas para mejorar la calidad del código y responder rápidamente a cambios en los requisitos.

\subsubsection{¿Para qué sirve?}
XP es ideal para proyectos donde la calidad del código es crítica, cambios frecuentes son esperados y se requiere reducir riesgos técnicos y defectos en etapas tempranas.

\subsubsection{¿Cómo se implementa?}
\begin{itemize}
    \item Pair Programming: programación en parejas
    \item Test-Driven Development (TDD): escribir pruebas antes del código
    \item Integración continua: fusionar código frecuentemente
    \item Refactorización continua: mejorar código sin cambiar funcionalidad
    \item Propiedad colectiva del código: cualquiera puede modificar
    \item Estándares de código estrictos
    \item Releases frecuentes
\end{itemize}

\section{Comparativa y Análisis}

\begin{table}[h]
\centering
\begin{tabular}{|l|l|l|l|}
\hline
\textbf{Metodología} & \textbf{Flexibilidad} & \textbf{Documentación} & \textbf{Velocidad} \\
\hline
Agile & Alta & Media & Alta \\
\hline
Waterfall & Baja & Alta & Baja \\
\hline
Lean & Media & Baja & Alta \\
\hline
Black Belt & Media & Alta & Media \\
\hline
\end{tabular}
\end{table}

\section{Conclusiones}

La elección de una metodología de manejo de proyectos depende de múltiples factores:
\begin{itemize}
    \item Naturaleza del proyecto
    \item Tamaño y experiencia del equipo
    \item Requisitos de cambio esperados
    \item Restricciones de tiempo y presupuesto
    \item Necesidades de documentación
\end{itemize}

Un proyecto de software exitoso requiere seleccionar la metodología más adecuada o hibridar varias metodologías según las necesidades específicas del contexto.

\end{document}
